% Created 2026-08-29 Sat 19:34
% Intended LaTeX compiler: pdflatex
\documentclass[11pt]{article}
\usepackage[utf8]{inputenc}
\usepackage[T1]{fontenc}
\usepackage{graphicx}
\usepackage{longtable}
\usepackage{wrapfig}
\usepackage{rotating}
\usepackage[normalem]{ulem}
\usepackage{amsmath}
\usepackage{amssymb}
\usepackage{capt-of}
\usepackage{hyperref}
\usepackage[margin=0.5in]{geometry}
\usepackage{parskip}
\usepackage{booktabs}
\usepackage{array}
\usepackage{ragged2e}
\usepackage{tabularx}
\usepackage{minted}
\usepackage{makeidx}
\newcommand{\maxtablewidth}{7in}
\newcolumntype{L}{>{\RaggedRight\arraybackslash}X}
\newcolumntype{C}{>{\Centering\arraybackslash}X}
\setlength{\emergencystretch}{3em}
\author{Blaine Mooers}
\date{\today}
\title{A PyQt5 Point-and-Click Front End for writing-habit and writing-schedule}
\hypersetup{
 pdfauthor={Blaine Mooers},
 pdftitle={A PyQt5 Point-and-Click Front End for writing-habit and writing-schedule},
 pdfkeywords={},
 pdfsubject={},
 pdfcreator={Emacs 30.2 (Org mode 9.7.11)}, 
 pdflang={English}}
\begin{document}

\maketitle
\section{Summary}
\label{sec:org37802ad}

The two packages already separate their work cleanly into a parsing layer, a
generation layer, and a thin \texttt{argparse} front end, so a graphical interface
adds a third front end rather than a rewrite. The plan below places that
interface in the \texttt{writing-habit} repository at \texttt{src/writing\_habit/gui/}, behind
an optional \texttt{[gui]} extra, and gives it two halves. The first half is a
spreadsheet-style editor for the weekly template table, which is the single
plain-text artifact both packages read. The second half is a command panel
whose forms are generated by reading the existing \texttt{argparse} parsers, so every
subcommand and every option becomes a widget without a hand-maintained
duplicate list.

One design decision carries most of the value. The command panel builds itself
from \texttt{writing\_habit.cli.build\_parser()} and \texttt{writing\_schedule.cli.build\_parser()}
rather than from a hard-coded catalog of commands, because a hand-written
catalog drifts the first time a flag is added. A second decision protects the
command line as the reference interface. Every action in the graphical
interface displays the exact shell command it is about to run, so the interface
teaches the command line rather than hiding it.

The path you named, \texttt{\textasciitilde{}/6112MooersLabGitHubLabRepos/writing-schedule-py/writing\_schedule},
holds the scheduler that owns the template table. The tracker that you named in
the request lives beside it at \texttt{\textasciitilde{}/6112MooersLabGitHubLabRepos/writing-habit}.
The plan covers both, because the template editor needs the scheduler parser
and the command panel needs the tracker commands.
\section{Decisions taken}
\label{sec:org1cb7106}
\index{decisions}

The ten open questions of the first draft are now settled, and the plan below
reflects the answers rather than the alternatives.

\begin{table}[htbp]
\caption{The decisions that fix the design, and the consequence of each.}
\centering
\begin{tabularx}{\maxtablewidth}{LL}
\toprule
Decision & Consequence\\
\midrule
PyQt5 is available to the interpreter you will use & no version negotiation is needed, and the pin stays at 5.15\\
The graphical layer ships inside \texttt{writing-habit} & one repository, one release, and an optional \texttt{[gui]} extra\\
PyQt5 under the GPL & the README states the split license, because the library and the command line remain MIT\\
A repeated day pattern that is not adjacent is written out again & the encoder collapses maximal runs of consecutive identical days only\\
A change of project repeats the activity letter, so \texttt{sBBBsCCCsD} stands & a run covers consecutive blocks sharing activity and project, and two template names need renaming\\
\texttt{:risky:} replaces \texttt{:speculative:} as the tag in a table & the stored class keeps the name \texttt{speculative}, so the schema, the views, both dashboards, the fixtures, and the elisp twin are untouched\\
\texttt{:speculative:} is no longer read & an old table is flagged as carrying an unread tag rather than misread, and the shipped tables were rewritten\\
\texttt{name.encode} joins \texttt{name.decode} & the editor can name the file it saves, and the command line gains the inverse function\\
A \texttt{dashboard} subcommand joins the command line & the two front ends offer the same actions\\
The table round-trip test comes before the editor writer & the one unrecoverable failure is caught first\\
\texttt{pytest-qt} joins the \texttt{dev} extra, and the Makefile exports \texttt{QT\_QPA\_PLATFORM=offscreen} & the suite runs without a display\\
A per-command database path stays & each command form carries its own \texttt{-{}-{}db} widget\\
The Sphinx site gains a page for the interface & the graphical layer is documented where the rest of the package is\\
\bottomrule
\end{tabularx}
\end{table}

Every step of the build order below is complete. The interface edits the weekly
table, runs every subcommand of both programs, shows what each run wrote, and is
documented on the Sphinx site. Step 0 is the library
work these decisions imply, step 1 is the window shell with its packaging, step
2 is the generated command panel that runs every tracker subcommand, and step 3
extends it to the scheduler. Every subcommand of both programs now runs by point
and click, which leaves the weekly-table editor of steps 4 through 6.
\section{What exists today}
\label{sec:org40427b2}

\subsection{The two command surfaces}
\label{sec:orgf1729a4}
\index{writing-habit command line}
\index{writing-schedule command line}

Every widget in the plan maps onto one of the subcommands below, which were
read from the two \texttt{cli.py} files rather than from the documentation.

\begin{table}[htbp]
\caption{The subcommands of the two packages and the options each accepts.}
\centering
\begin{tabularx}{\maxtablewidth}{LLL}
\toprule
Package & Subcommand & Options\\
\midrule
writing-habit & \texttt{initdb} & \texttt{-{}-{}db}\\
writing-habit & \texttt{plan import PATH} & \texttt{-{}-{}week}, \texttt{-{}-{}db}\\
writing-habit & \texttt{track import PATH} & \texttt{-{}-{}format \{csv,ics\}}, \texttt{-{}-{}db}\\
writing-habit & \texttt{track add} & \texttt{-{}-{}day}, \texttt{-{}-{}project}, \texttt{-{}-{}minutes}, \texttt{-{}-{}category}, \texttt{-{}-{}start}, \texttt{-{}-{}end}, \texttt{-{}-{}note}, \texttt{-{}-{}db}\\
writing-habit & \texttt{compare} & \texttt{-{}-{}week}, \texttt{-{}-{}plot}, \texttt{-{}-{}db}\\
writing-habit & \texttt{dashboard} & \texttt{-{}-{}week}, \texttt{-{}-{}out}, \texttt{-{}-{}db}\\
writing-habit & \texttt{history} & \texttt{-{}-{}from}, \texttt{-{}-{}to}, \texttt{-{}-{}plot}, \texttt{-{}-{}db}\\
writing-habit & \texttt{context set} & \texttt{-{}-{}week}, \texttt{-{}-{}tag}, \texttt{-{}-{}note}, \texttt{-{}-{}db}\\
writing-habit & \texttt{context clear} & \texttt{-{}-{}week}, \texttt{-{}-{}tag}, \texttt{-{}-{}db}\\
writing-habit & \texttt{context list} & \texttt{-{}-{}week}, \texttt{-{}-{}db}\\
writing-habit & \texttt{seasons} & \texttt{-{}-{}out}, \texttt{-{}-{}db}\\
writing-habit & \texttt{name CODE} & \texttt{-{}-{}table}\\
writing-schedule & \texttt{generate TABLE} & \texttt{-{}-{}week}, \texttt{-{}-{}day}, \texttt{-{}-{}dir}, \texttt{-{}-{}no-ics}, \texttt{-{}-{}no-todo}, \texttt{-{}-{}tz}, \texttt{-{}-{}strict}\\
writing-schedule & \texttt{export TABLE} & \texttt{-{}-{}week}, \texttt{-{}-{}day}, \texttt{-{}-{}dir}, \texttt{-{}-{}no-todo}, \texttt{-{}-{}tz}, \texttt{-{}-{}strict}\\
writing-schedule & \texttt{sheets TABLE} & \texttt{-{}-{}week}, \texttt{-{}-{}day}, \texttt{-{}-{}dir}, \texttt{-{}-{}engine \{reportlab,latex\}}, \texttt{-{}-{}format \{pdf,org,both\}}, \texttt{-{}-{}per-day}, \texttt{-{}-{}strict}\\
writing-schedule & \texttt{check TABLE} & none\\
writing-schedule & \texttt{template N} & \texttt{-{}-{}out}\\
writing-schedule & \texttt{weeks} & \texttt{-{}-{}dir}\\
\bottomrule
\end{tabularx}
\end{table}

Two observations came out of that reading. The \texttt{dashboard} subcommand described
in the project status note was absent from \texttt{writing\_habit/cli.py}, although
\texttt{dashboard.write\_dashboard(con, week, out\_path)} existed and the golden fixture
test covered it. Step 0 added \texttt{writing-habit dashboard -{}-{}week WEEK -{}-{}out PATH
-{}-{}db DB}, so the table above lists it and the two front ends offer the same set
of actions. The \texttt{writing-schedule} package also carries a
\texttt{check} subcommand that the status notes do not mention, and it is the natural
backing call for the live overlap indicator in the editor.
\subsection{The template table}
\label{sec:orgf559424}
\index{weekly block table}
\index{org table format}

The template is a single org table. Its first row names the days, its section
rows name the activity class, its time-block rows carry one project code per
day cell, and its legend rows map each code to a description and an optional
risk class. The parser in \texttt{writing\_schedule/parser.py} tests four row kinds in a
fixed order and takes the first match, so the editor model must preserve that
row order when it writes the file back.

\begin{table}[htbp]
\caption{The four row kinds the parser recognizes, and how the editor represents each.}
\centering
\begin{tabularx}{\maxtablewidth}{LLL}
\toprule
Row kind & Recognized by & Editor representation\\
\midrule
Header & any cell after the first naming a day, such as \texttt{M} or \texttt{Tue} & column headers of the grid, fixed\\
Section & first cell of letters and spaces, optional trailing colon & a spanning group row, not editable as a block\\
Time block & a time range such as \texttt{04:00-05:30} in the first cell & an editable row of one combo box per day\\
Legend & first cell matching a short uppercase code then a colon & a row in the separate legend table\\
\bottomrule
\end{tabularx}
\end{table}

The risk class rides in the legend description as either \texttt{(safe)} or \texttt{:safe:},
and \texttt{writing\_habit/name.py} strips it when reading the legend. The editor
therefore edits the risk class as a combo box and writes the tag back in the
form the file already used, so a round trip does not rewrite unrelated bytes.
\subsection{What the tracker adds}
\label{sec:org6f287bd}
\index{schedule file-name code}
\index{adherence database}

The tracker reads the same table through \texttt{plan\_import.import\_org}, stores the
schedule code taken from the file name, and keeps every planned block and every
tracked session in one SQLite database. \texttt{name.decode} turns a code such as
\texttt{4gAAeAsA-gWW} into a week of blocks, \texttt{name.summary} counts them, and
\texttt{name.check\_against\_legend} reports whether each project letter resolves to a
legend entry. Those three functions give the editor a live readout of what the
current grid means, at no cost beyond calling them.
\section{Architecture}
\label{sec:org3fc4470}

\subsection{In-process library calls, with the shell command on display}
\label{sec:org711a2d8}
\index{architecture decision}

Two ways of driving the tools are possible. The interface can shell out to the
\texttt{writing-habit} and \texttt{writing-schedule} executables, or it can import the
packages and call their functions. The plan calls the library functions in a
worker thread, for four reasons. The functions return Python objects, so the
grid, the overlap list, and the report tables need no output parsing. Errors
arrive as exceptions with tracebacks rather than as exit codes. Startup avoids
a subprocess for every keystroke-driven revalidation. Testing needs no
executable on the path.

The cost of that choice is that the graphical interface can drift away from the
command line. Build step 2 removed the cost rather than managing it. The worker
calls \texttt{writing\_habit.cli.main(argv)} with the argument list the form built, so
the preview above the run button is not merely equivalent to what runs; it is
what runs. Drift is impossible while that holds, and a test asserts the stronger
property, namely that the generated argument list parses back through the
parser to the values the form holds. Every run records the command string in a
session log pane the user can copy.

The interface drives two programs, so the runner keeps a table naming the module
that holds each \texttt{main} function and imports it on demand. Resolving the program
by name rather than by an import at load time is what lets the tracker run when
the optional scheduler is absent. A tab whose commands belong to a missing
program shows the install line rather than vanishing, and the status bar says
which program is missing, because a silently shorter window teaches the user
nothing.

Two consequences of running in process deserve a note. The output capture
replaces \texttt{sys.stdout} for the whole process, so a module lock allows one command
at a time and the panel disables its run button while a command is in flight.
The library uses \texttt{sys.exit} with a message to explain a missing optional
dependency, so the runner turns a string payload into exit code one and writes
the message to the captured output rather than discarding it. A user without
\texttt{writing-schedule} installed therefore reads the install hint in the output pane
instead of a traceback.
\subsection{Module layout}
\label{sec:org17500ff}
\index{module layout}

\begin{table}[htbp]
\caption{The proposed modules under \texttt{src/writing\_habit/gui/} and the responsibility of each.}
\centering
\begin{tabularx}{\maxtablewidth}{LL}
\toprule
Module & Responsibility\\
\midrule
\texttt{qt.py} & the single import site for Qt, so a later move to PyQt6 or PySide6 touches one file\\
\texttt{app.py} & the \texttt{main()} entry point, the \texttt{QApplication}, and the top-level exception hook\\
\texttt{mainwindow.py} & the window, the tab stack, the menu bar, the status bar, and the session log dock\\
\texttt{settings.py} & \texttt{QSettings} for the database path, the table directory, the output directory, and the recent files\\
\texttt{argspec.py} & introspection of an \texttt{argparse} parser into a list of command and option descriptors\\
\texttt{forms.py} & one option descriptor to one widget, plus the reverse read of the widget into an argument list\\
\texttt{runner.py} & a \texttt{QThread} worker that runs one command, captures \texttt{stdout} and \texttt{stderr}, and emits result or error\\
\texttt{command\_panel.py} & the generated form for one subcommand, its command preview, its run button, and its output pane\\
\texttt{weekly\_table.py} & the table as a document, classified by row kind, with the derived readings and a byte-exact writer. No Qt\\
\texttt{table\_model.py} & a \texttt{QAbstractTableModel} over that document, mapping a grid cell back to its row and cell\\
\texttt{schedule\_editor.py} & the grid, the four reading panels, and the open and reload controls\\
\texttt{legend\_model.py} & a \texttt{QAbstractTableModel} over the legend rows, with code, description, and risk class\\
\texttt{report\_view.py} & \texttt{ArtifactView}, which renders one file a command wrote, and \texttt{PreviewDock}, which keeps a tab per recent file\\
\bottomrule
\end{tabularx}
\end{table}
\subsection{Tab layout}
\label{sec:org97cc24d}
\index{window layout}

The window carries one tab per stage of the weekly loop, in the order the loop
runs.

\begin{table}[htbp]
\caption{The tabs, the widgets each holds, and the calls each makes.}
\centering
\begin{tabularx}{\maxtablewidth}{LLL}
\toprule
Tab & Widgets & Backing calls\\
\midrule
Schedule & the weekly grid and the legend table, above the generated \texttt{template} and \texttt{check} forms & \texttt{parse\_text}, \texttt{find\_overlaps}, \texttt{name.decode}, \texttt{name.encode}, and \texttt{writing-schedule} \texttt{template} and \texttt{check}\\
Generate & week or day picker, output directory, the \texttt{-{}-{}no-ics}, \texttt{-{}-{}no-todo}, \texttt{-{}-{}tz}, and \texttt{-{}-{}strict} controls & \texttt{writing-schedule} \texttt{generate}, \texttt{export}, and \texttt{weeks}\\
Sheets & engine and format combo boxes, the per-day check box, a page preview & \texttt{writing-schedule} \texttt{sheets}\\
Plan & table picker, week picker, database picker & \texttt{plan\_import.import\_org}\\
Track & a file importer and a hand-entry form for one session & \texttt{track.csv\_actuals.import\_csv}, \texttt{track.ics\_actuals.import\_ics}, \texttt{track.manual.add\_session}\\
Compare & week picker, the text report, the optional plot, the dashboard preview & \texttt{compare.report.render\_week}, \texttt{report.write\_plot}, \texttt{dashboard.write\_dashboard}\\
History & a from and to range, the text series, the plot & \texttt{compare.history.render\_text}, \texttt{history.write\_plots}\\
Context & a week picker, a tag combo box, a note field, the tag list & \texttt{context.set\_tag}, \texttt{clear\_tag}, \texttt{list\_tags}\\
Seasons & an output path and a preview & \texttt{seasons.write\_seasons}\\
\bottomrule
\end{tabularx}
\end{table}
\section{The schedule editor}
\label{sec:org00d8c71}

\subsection{The grid}
\label{sec:org3840fd8}
\index{schedule grid}
\index{combo box delegate}

The grid is a \texttt{QTableView} over \texttt{table\_model.py}, which reads
\texttt{weekly\_table.py}. Columns are the days found in the header row, so a six-day
table and a seven-day table both work without a setting. Rows follow the file
order, which keeps section rows in place. A section row renders bold and shaded
and refuses edits.

The document underneath is the part that matters. It keeps every line of the
file, classified by the four row kinds in the order the scheduler parser tests
them, because that parser takes the first match and a different order here would
give the editor a different table from the one the tools read. Two properties
are asserted over every shipped table rather than assumed. The document sees the
same blocks the scheduler sees, and an untouched document writes back byte for
byte. The second property is the round-trip test the build order puts before the
writer, and it holds for all thirteen tables in \texttt{templates/} and \texttt{examples/}.

A time-block cell is edited through a \texttt{QStyledItemDelegate} presenting a combo
box of the legend codes plus a blank entry, so a cell rarely holds a code the
legend does not define, though the box stays editable for a writer who wants to
type a code and add its legend row after. A section row and the time column
refuse an edit, because they belong to the file rather than to the week.

The legend has its own table below the grid, with the code, the description, and
the risk tag as a choice of \texttt{none}, \texttt{safe}, and \texttt{risky}. The choice shows the
tag a writer types rather than the class the database stores, because the table
is the thing being edited. Writing \texttt{risky} puts \texttt{:risky:} in the file and yields
the class the database calls \texttt{speculative}.

Every edit rewrites exactly one line. A value that fits its column keeps the
table aligned, and a longer value widens its own slot alone, which leaves the
table misaligned until the writer realigns it in Emacs. Reflowing lines they did
not touch would be worse, because the diff would then hide the change they made.

Three edits change the shape of the week rather than one cell. Adding a time
block inserts a row under the current section and asks for the time range with
a validated line edit. Adding a section appends a section row and one block
row. Adding a day is out of scope for the first release, because the day set
comes from the header row and the parser reads it once.
\subsection{The live readouts}
\label{sec:org89b147a}
\index{live validation}
\index{overlap warning}

One revalidation pass runs when a table opens, and step 5 will run the same pass
after every edit, which is cheap because a week holds a few dozen blocks. The
pass fills four panels beside the grid.

The name panel gives the canonical file name for the grid, the week the code
decodes to, and the block counts by activity and by project. When the file name
on disk differs from the canonical form, the panel says so, which is how a
writer notices that a table has drifted from its name.

The clashes panel lists the pairs from \texttt{writing\_schedule.overlap.find\_overlaps},
the tab carries the count, and the offending cells take a warning tint. This is
the graphical equivalent of the \texttt{check} subcommand and of the \texttt{-{}-{}strict} guard.

The totals panel shows planned minutes by day, by project, and by activity, with
the week total. It also names any section header that maps to no activity
category, because the plan importer silently counts such a section as
generative, and a writer should know that before the numbers reach the database.

The legend panel shows the output of \texttt{name.check\_against\_legend}, so an unknown
or ambiguous project letter is visible before the file is used. It also lists
legend entries the week never uses, the codes defined more than once, and the
risk tags that name no class.

One reading has no answer for some tables, and saying so plainly is the useful
behavior. A file name gives one letter to one block, so a legend code of more
than one letter, such as \texttt{EM} for email, cannot appear in a cell of a named
table. The name panel then says which code is at fault and that the naming rules
answer it with a single-letter alias, because only the writer knows which letter
is free.
\subsection{What the real week showed}
\label{sec:org1591c9c}
\index{real-world table}

\texttt{examples/aug24.org} is a real weekly table rather than a demonstration, so
running the document over it was the first honest test. The table offers a menu
of one hundred and twenty-six candidate time blocks and fills five of them, all
on the Thursday, which the readings handle correctly. The week totals five
hundred and seventy planned minutes, nothing clashes because the blocks merely
touch, and the canonical name is \texttt{3o-gHgAeCsIsB}, where the three open days are
Monday through Wednesday. An empty cell is not a plan, and only a lettered cell
becomes a block.

Three faults in the data surfaced that no synthetic example would have shown.

\begin{table}[htbp]
\caption{What the real table revealed, and what each fault costs.}
\centering
\begin{tabularx}{\maxtablewidth}{LLL}
\toprule
Fault & Cost & Response\\
\midrule
The legend defines \texttt{H} twice, for \texttt{1006AIrxOpt} and \texttt{2174ComputationalCrystallography}, and skips \texttt{O} & one project is silently unreachable, and the Thursday \texttt{H} block resolves to the first & the legend panel names the code and both descriptions; the writer chooses\\
\texttt{G} and \texttt{P} carried \texttt{:spec:}, which names no risk class & both projects fell out of the barbell comparison, which reads only the tags it knows & the legend panel reports a tag that was not read; the table now reads \texttt{:risky:}\\
\texttt{name.read\_legend} kept the last definition of a duplicated code while the plan importer keeps the first & one table gave two different descriptions for one project, depending on the command & \texttt{read\_legend} now keeps the first, matching the importer, with a test asserting the agreement\\
\bottomrule
\end{tabularx}
\end{table}

The third fault was in the code rather than the data, and it existed before this
work. The scheduler resolves a legend code in the manner of elisp \texttt{assoc}, which
returns the first matching entry, so the importer and the reader disagreed on
exactly the tables where it matters.
\subsection{The rounding that made a page depend on its machine}
\label{sec:org82a33f5}
\index{portable rounding}

\texttt{tests/test\_seasons.py} failed from the first day of this work, and the cause
turned out to be worth more than the test. The grouping views ended with
\texttt{ROUND(1.0*SUM(actual\_min)/SUM(planned\_min), 2)}, and both dashboards printed
that column. One row of the fixture is 510 minutes against 1200, which is
exactly 0.425, and SQLite does not agree with itself across versions on such a
value. Version 3.37 gives 0.43 and version 3.45 gives 0.42, so one database
rendered two different pages on two machines, and the golden that both ports
assert against failed on whichever machine had not written it.

The fix is to compute every displayed ratio where the page is built rather than
in the query. Python formats the same double the same way everywhere, and so
does the elisp twin, so the page a database yields no longer depends on the
SQLite build underneath. The views keep their \texttt{adherence} column, because
nothing forces a stored view to change and the fixture databases carry the old
definition anyway.

The committed goldens did not move, which is the useful part. The computed value
reproduces both pages byte for byte, so the fix removed a dependency rather than
redefining the numbers. Two tests hold the line: one asserts that a view column
which disagrees never reaches the page, and one pins the ratio of the value the
two SQLite builds disputed.
\subsection{The encoder}
\label{sec:org30ee4eb}
\index{name encoder}
\index{canonical file name}

\texttt{name.py} decoded a code into a week and had no inverse, so the editor had no
way to name the file it saves. \texttt{name.encode(week)} and \texttt{name.encode\_day(blocks)}
now supply that inverse. They live in the library rather than in the graphical
layer, because the command line and the Emacs Lisp twin want the same function.

The canonical spelling follows \texttt{docs/table-file-naming-rules.org}. A maximal run
of identical consecutive days collapses into one group carrying the day count, a
count of one is omitted, trailing open days are dropped because they are
implied, and an interior open day is written as \texttt{o}. Two days that share a
pattern without being adjacent are written out again, which is the rule you
chose, so a Monday, Wednesday, Friday week reads \texttt{gA-o-gA-o-gA}. The property
\texttt{encode(decode(code)) =} code= holds for every canonical code, and the corpus
test asserts the weaker semantic property \texttt{decode(encode(decode(name))) =}
decode(name)= for every file name in \texttt{templates/}.

A second rule was needed for a run of blocks that share an activity while the
project changes. Both spellings decode to the same week, so the six canonical
rules did not choose between them, and the corpus splits evenly. Two template
names repeat the activity letter at each change of project, giving
\texttt{sBBBsCCCsD}, and two merge the run, giving \texttt{gAAB}. The remaining eight names
never change project inside a run, so they are silent on the question.

The repeated activity letter is the canonical form, which is your decision. A
run therefore covers consecutive blocks that share both the activity and the
project, so three support blocks on B followed by three on C followed by one on
D read \texttt{sBBBsCCCsD} rather than \texttt{sBBBCCCD}. The rule keeps every project visible
in the file name, which is the point of the name. The two names that use the
merged spelling inside a generative run become non-canonical, so
\texttt{gAABeAsA-gAABeBsB-gBBAeAsB-gBBAeBsW-gWsA.org} and its six-day sibling would be
renamed to \texttt{gAAgBeAsA-gAAgBeBsB-gBBgAeAsB-gBBgAeBsW-gWsA.org}. Nothing renames
them automatically, because the tracker groups weeks by the stored code.
\subsection{Saving}
\label{sec:orgfd7669a}
\index{file naming}
\index{round trip}

Saving writes the document back to the file it came from, and \texttt{Save as} offers
the canonical name \texttt{<code>.org} when the grid has one. A file that was not edited
saves byte for byte, and a file that was edited differs only on the lines that
changed. Those two properties were tested before the writer existed, because
they are the failure a writer cannot recover from.

\texttt{New from template} asks for a project count and calls
\texttt{writing\_schedule.template.template\_string}, so a blank week comes from the
scheduler's own scaffold rather than from a second copy of the layout here.

\texttt{Rename to canonical} moves the file to the name its grid gives it. The button
offers itself only when the file is saved and misnamed, and the rename refuses
to overwrite a different file, refuses a week that has no canonical name, and
refuses while edits are unsaved. The action matters more than it looks, because
the tracker groups weeks by the schedule code stored at plan import, and that
code comes from the file name, so a drifted name files the week under the wrong
plan shape and the seasonal comparison then groups it with the wrong weeks.
\section{The command panel}
\label{sec:org21fa515}

\subsection{Generating forms from argparse}
\label{sec:org741872b}
\index{argparse introspection}
\index{generated forms}

\texttt{argspec.py} walks a parser and yields a descriptor for every subcommand and
every option. The walk reads \texttt{parser.\_subparsers}, then for each subparser reads
\texttt{\_actions}, skipping the help action. Each action gives the option strings, the
\texttt{dest}, the \texttt{help} text, the \texttt{choices}, the \texttt{type}, the \texttt{required} flag, and
whether it is a \texttt{store\_true} action. Private attributes are involved, so the
walk is confined to \texttt{argspec.py} and is covered by a test that fails loudly when
a future Python changes the layout.

\begin{table}[htbp]
\caption{The rule that maps an argparse action to a widget.}
\centering
\begin{tabularx}{\maxtablewidth}{LLL}
\toprule
Action shape & Widget & Example\\
\midrule
\texttt{action}"store\textsubscript{true}"= & \texttt{QCheckBox} & \texttt{-{}-{}strict}, \texttt{-{}-{}no-ics}, \texttt{-{}-{}per-day}\\
\texttt{choices} present & \texttt{QComboBox} & \texttt{-{}-{}format}, \texttt{-{}-{}engine}\\
\texttt{type=int} & \texttt{QSpinBox} & \texttt{-{}-{}minutes}\\
\texttt{dest} in the date set & \texttt{QDateEdit} with a calendar popup and a "today" button & \texttt{-{}-{}week}, \texttt{-{}-{}day}, \texttt{-{}-{}from}, \texttt{-{}-{}to}\\
\texttt{dest} naming a directory & line edit plus a directory chooser & \texttt{-{}-{}dir}\\
\texttt{dest} naming a file or a positional path & line edit plus a file chooser with a suffix filter & \texttt{-{}-{}out}, \texttt{-{}-{}plot}, \texttt{-{}-{}table}, \texttt{path}\\
anything else & \texttt{QLineEdit} & \texttt{-{}-{}tag}, \texttt{-{}-{}note}, \texttt{-{}-{}tz}\\
\bottomrule
\end{tabularx}
\end{table}

The sets are small tables in \texttt{forms.py}, and they are keyed by the primary flag
rather than by \texttt{dest}, because two commands can share a \texttt{dest} while meaning
different things. \texttt{history -{}-{}from} and \texttt{track add -{}-{}start} both carry the \texttt{dest}
\texttt{start}, yet one is a calendar date and the other is a clock time. A test
asserts that pair, because it is the case a \texttt{dest} table would silently get
wrong.

One more table holds suggestions. \texttt{-{}-{}category} and \texttt{-{}-{}tag} take a value the
parser does not restrict, while the schema and the documented vocabulary do
have expected values, so those two get an editable combination box that offers
\texttt{generative}, \texttt{editing}, and \texttt{support} for the first and the event vocabulary
for the second. A new option that falls outside every table still produces a
working line edit, so an unrecognized flag degrades to a usable widget rather
than to a crash.

An optional value needs a way to say nothing at all. An empty line edit omits
its option, and a date, a time, or a count carries a small check box that
decides whether the option is sent, because those widgets always show
something.

One option needs the opposite. \texttt{writing-schedule -{}-{}tz ""} asks for floating
calendar times, so an empty value there is a request rather than an omission.
That flag gets a line edit with the same send box, and the argument builder
treats ``None`` as the omission and an empty string as a value.
\subsection{Running a command}
\label{sec:orgf1c654e}
\index{worker thread}
\index{sqlite thread safety}

The run button collects the widget values into an argument list, renders the
preview string, and hands a callable to \texttt{runner.py}. The worker runs on a
\texttt{QThread} and communicates only through signals, because touching a widget from
a worker thread is the standard way to crash a Qt program. The worker
redirects \texttt{sys.stdout} and \texttt{sys.stderr} with \texttt{contextlib.redirect\_stdout} so that
the printed output of the existing functions reaches the log pane unchanged.

The database needs care. \texttt{sqlite3} connections belong to the thread that
created them, so the worker calls \texttt{db.connect} itself and closes the connection
when the task ends. The main thread keeps no long-lived connection. Read-only
panels that need a quick query open a short connection on the main thread and
close it immediately, which is acceptable at this data volume.
\subsection{Preferences and defaults}
\label{sec:org048b0fd}
\index{QSettings}
\index{default paths}

The database path, the template directory, the output directory, the time zone,
and the last used week persist through \texttt{QSettings}. Every form seeds its widgets
from those settings, so the common case of running \texttt{plan import}, then
\texttt{compare}, then \texttt{dashboard} for the same week needs no retyping.

The \texttt{-{}-{}db} widget stays on every command form, which is your decision and which
matches the command line, where \texttt{-{}-{}db} is required on every subcommand that
touches the database. Two habits keep a per-command path from becoming a source
of silent mistakes. The remembered path seeds every form, so the widgets agree
unless you deliberately change one. The status bar shows the path of the
database the last command actually used, together with its row counts, so a
command that ran against the wrong file is visible immediately rather than three
commands later.
\section{Seeing what a command wrote}
\label{sec:orgceefa4a}
\index{preview dock}
\index{web engine}

A command that writes a file used to end with a line of text naming a path, and
the writer then went looking for it. The preview dock removes that step. It
holds one tab per recent file, newest last, and it appears the first time a
command produces something rather than occupying the window before then.

The panel knows what its run produced in two ways, and it needs both. The paths
given as \texttt{-{}-{}out} and \texttt{-{}-{}plot} come from the form, so nothing has to be parsed
for the tracker commands. The scheduler instead takes a directory and derives
the file names, so those paths are read from the \texttt{Wrote} lines the command
prints, taking the last token of each line and keeping it only when it names a
file on disk. A token that is not a file is dropped, so a sentence that merely
begins with the word is harmless.

\begin{table}[htbp]
\caption{What each kind of output does in the preview.}
\centering
\begin{tabularx}{\maxtablewidth}{LLL}
\toprule
Output & Shown as & Note\\
\midrule
The weekly dashboard and the seasons page & the web engine when it is installed, and the rich-text widget otherwise & the fallback says what it cannot render\\
The compare and history plots & the image, with its pixel size & scrolls when it is larger than the dock\\
The schedule \texttt{.org}, the \texttt{.ics}, the sheet \texttt{.org}, and the \texttt{.tex} & plain text & a file over two hundred kilobytes is named rather than loaded\\
The sheet PDFs & the name and nothing else & Qt has no PDF viewer here, so the desktop opens it\\
\bottomrule
\end{tabularx}
\end{table}

Every kind offers \texttt{Open outside}, which hands the file to the program the
desktop would use. That button is the honest answer for the dashboards, because
they use CSS that the rich-text widget of Qt renders only roughly.
\subsection{The rule the web engine imposes}
\label{sec:orgcaeb8e3}
\index{PyQtWebEngine}

\texttt{PyQtWebEngine} ships separately from PyQt5, so most installations lack it and
the fallback is the normal path. Where it is installed, Qt imposes an order that
is easy to get wrong and quiet when it is. The web-engine module must be
imported, and the shared OpenGL context attribute must be set, before the
\texttt{QApplication} exists. Import it later and Qt prints a warning and refuses,
which leaves a machine that has the faithful renderer using the rough one
forever, with nothing in the interface to say why.

The startup therefore calls one preparation function before building the
application object, and the probe caches its answer so a later call still gives
the right one. A test asserts the order inside \texttt{main}, because the failure is
invisible at run time.
\section{Packaging}
\label{sec:orge5c41c7}

\subsection{The optional extra}
\label{sec:org695e99f}
\index{optional dependency}
\index{entry point}

The graphical layer ships inside \texttt{writing-habit} behind an extra, so the wheel
and the command line stay free of a Qt dependency.

\begin{verbatim}
[project.optional-dependencies]
gui = ["PyQt5>=5.15", "writing-schedule>=0.3.1"]
preview = ["PyQtWebEngine>=5.15"]

[project.scripts]
writing-habit = "writing_habit.cli:main"
writing-habit-gui = "writing_habit.gui.app:main"
\end{verbatim}

The web engine sits in its own extra rather than in \texttt{gui}, because \texttt{Open
outside} already answers the same need and the engine is a large dependency for
a convenience. It is also GPL, so an installation that wants neither GPL package
simply installs neither extra.

Three rules keep the extra honest. No module outside \texttt{writing\_habit/gui/} may
import Qt. The \texttt{gui} package imports Qt only through \texttt{gui/qt.py}. The
\texttt{writing-habit-gui} entry point catches the import error and prints the
installation hint \texttt{pip install "writing-habit[gui]"} rather than a traceback.

The \texttt{gui} extra depends on \texttt{writing-schedule} because the editor uses its parser,
its overlap detector, and its template scaffold. That package is already an
optional dependency under the \texttt{plan} extra, so the two extras should be
documented together.
\subsection{The choice of PyQt5 and the split license}
\label{sec:orgc71b6ec}
\index{PyQt5 versus PyQt6}
\index{licensing}

PyQt5 is the chosen binding, and it is available to the interpreter you will
use. PyQt5 is in maintenance while PyQt6 and PySide6 receive the new work, so
the single import site in \texttt{gui/qt.py} earns its keep as insurance. A later move
changes the import names and a small number of enumeration references, which is
a one-file change when nothing else imports Qt directly.

PyQt5 is distributed under the GPL, so the package now carries two licenses. The
\texttt{pyproject.toml} comment above the extra records that, and three places need the
same sentence. The README needs a license section that states that the library
and the command line are MIT while the optional graphical extra pulls in PyQt5
under the GPL. The \texttt{About} dialog of the interface needs the same statement. The
Sphinx installation page needs it beside the \texttt{pip install "writing-habit[gui]"}
line. A user who installs the package without the extra is unaffected, which is
the reason the extra exists.
\section{Testing}
\label{sec:orgd5f8fa7}
\index{pytest-qt}
\index{headless testing}

The test suite runs headless with \texttt{QT\_QPA\_PLATFORM=offscreen}, and \texttt{pytest-qt}
joins the \texttt{dev} extra. Six groups of tests carry the weight.

\begin{table}[htbp]
\caption{The test groups, what each protects, and the failure each prevents.}
\centering
\begin{tabularx}{\maxtablewidth}{LLL}
\toprule
Group & What it asserts & Failure it prevents\\
\midrule
Round trip & every file in \texttt{templates/} and \texttt{examples/} loads into the model and writes back byte-identical & silent rewriting of a user's table\\
Model editing & a cell edit, a row insertion, and a legend edit produce the expected table text & corruption of row order\\
Encoder & \texttt{encode(decode(code)) =} code= for canonical codes, and the semantic round trip for every template name & a wrong file name, which becomes a wrong grouping key\\
Argument specification & every subcommand and option in both parsers produces a descriptor and a widget & a new flag that the interface cannot reach\\
Command preview & the rendered string parses back through the parser to the same namespace & a preview that lies about what will run\\
Worker & a failing task emits the error signal and leaves the window responsive & a frozen window on a bad path\\
\bottomrule
\end{tabularx}
\end{table}

The round-trip group is worth writing first, because it also documents the
table format in executable form.
\section{Build order}
\label{sec:orgb858123}
\index{milestones}

Each step below leaves a working program, so the work can stop at any step and
still be useful. Every step is complete.

\begin{enumerate}
\item {[}DONE] The library work the graphical layer depends on, none of which imports
Qt. \texttt{name.encode} and \texttt{name.encode\_day} with their tests, the \texttt{dashboard}
subcommand with its test, the \texttt{gui} extra and \texttt{pytest-qt} in the \texttt{dev} extra,
and a repository Makefile whose test target exports
\texttt{QT\_QPA\_PLATFORM=offscreen}.
\item {[}DONE] Add \texttt{gui/qt.py}, \texttt{gui/settings.py}, \texttt{gui/app.py}, and
\texttt{gui/mainwindow.py} with a tab per stage of the loop, the session-log dock,
the \texttt{About} dialog carrying the split license, and the \texttt{writing-habit-gui}
entry point. Two test modules come with it, one exercising the window
offscreen and one asserting that Qt never leaves the \texttt{gui} package. This
proves the packaging and the import guard before any feature exists.
\item {[}DONE] Build \texttt{argspec.py} and \texttt{forms.py} against
\texttt{writing\_habit.cli.build\_parser()}, and render the Plan, Track, Compare,
History, Context, and Seasons tabs from the descriptors. Add \texttt{runner.py} and
wire the panels to the log pane. All twelve tracker subcommands now run by
point and click, and the status bar names the database the last run used.
\item {[}DONE] Extend the same machinery to \texttt{writing\_schedule.cli.build\_parser()} for
the Schedule, Generate, and Sheets tabs. The runner gained a table of the two
programs and resolves the \texttt{main} function by name, so the tracker keeps
working when the scheduler is absent and those tabs then explain how to
install it. All six scheduler subcommands run by point and click.
\item {[}DONE] Build the read-only schedule grid, with the name, clashes, totals, and
legend panels. The document in \texttt{weekly\_table.py} carries the round-trip
property the writer of step 5 needs, and \texttt{name.parse\_legend\_cell} now serves
both the file reader and the in-memory grid, so one rule reads a legend row.
This step shows a table and writes nothing.
\item {[}DONE] Make the grid editable with the combo-box delegate, add the legend
editor with its risk-tag choice, and add the writer that rewrites one line
per edit. Add \texttt{Save}, \texttt{Save as} with the canonical name, and \texttt{New from
   template}, which calls \texttt{template\_string}.
\item {[}DONE] Wire the canonical file name into the interface and into the twin.
\texttt{name.encode} and its property test landed in step 0, and \texttt{Save as} offered
the canonical name in step 5, so this step closed the two gaps that left.
The editor gained a rename action, because a name that has drifted files the
week under the wrong shape and nothing acted on the panel's warning. The
Emacs Lisp twin gained \texttt{writing-habit-name-encode} and
\texttt{writing-habit-name-encode-day}, so both ports can name a week, and a
cross-check confirms they name every week alike.
\item {[}DONE] Add the preview dock for the dashboards, the plots, and the text
files, with \texttt{Open outside} for everything and the web engine where it is
installed.
\item {[}DONE] Add \texttt{docs/gui.md} to the Sphinx site, place it in the \texttt{index.md} table
of contents beside \texttt{cli.md}, and illustrate it with two screenshots taken
from the running interface. State the split license there, on the
installation page, and in the README. Add the \texttt{preview} extra for
\texttt{PyQtWebEngine}, which stays optional even with \texttt{gui}.
\end{enumerate}
\section{Risks and open questions}
\label{sec:org556f1fb}
\index{risks}

\begin{table}[htbp]
\caption{The risks that deserve a decision before coding, and the response to each.}
\centering
\begin{tabularx}{\maxtablewidth}{LLL}
\toprule
Risk & Why it matters & Response\\
\midrule
Private \texttt{argparse} attributes & \texttt{\_actions} and \texttt{\_subparsers} are not public API & confine the access to \texttt{argspec.py} and test it explicitly\\
The repeated activity letter & \texttt{sBBBsCCCsD} and \texttt{sBBBCCCD} decode alike, and the six canonical rules did not choose & the repeated letter is canonical; rename the two merged template names\\
Byte-exact round trip & column padding and trailing spaces vary between hand-edited files & normalize on load, and warn the user when the first save will reformat\\
GPL of PyQt5 beside MIT & a split license needs a plain statement & state it in the README, the \texttt{About} dialog, and the installation page\\
Threading with sqlite3 & a connection shared across threads raises at runtime & open and close the connection inside the worker\\
A per-command database path & one form can point at a different file than the last & seed every form from the remembered path, and name the database of the last run in the status bar\\
Feature drift from the command line & a graphical action with no command-line equivalent splits the two ports & require a subcommand for every action, which step 0 restored for the dashboard\\
The Emacs Lisp twin & the elisp package has no matching interface & treat the graphical layer as Python-only, and keep every shared rule in the library; the encoder and the \texttt{:risky:} tag are ported, and the two ports name every week alike\\
A legend code defined twice & the two readers disagreed, and one project becomes unreachable & \texttt{read\_legend} now keeps the first, and the legend panel reports the duplicate\\
An abbreviated risk tag & \texttt{:spec:} names no class, so the project leaves the barbell comparison without a word & the legend panel reports a tag that was not read\\
The web-engine import order & importing it after the \texttt{QApplication} exists leaves the faithful renderer unused and says nothing & one preparation call at startup, a cached probe, and a test on the order inside \texttt{main}\\
Rounding in the golden fixtures & \texttt{ROUND} in SQL rounded 0.425 differently between SQLite builds, so one database gave two pages & fixed: every displayed ratio is computed in Python, in both ports, and the goldens stand unchanged\\
\bottomrule
\end{tabularx}
\end{table}
\section{Follow-up checklist}
\label{sec:org58c37a5}

\begin{verbatim}
- [X] Confirm that PyQt5 is available to the Python interpreter you will use.
- [X] Ship the graphical layer inside writing-habit behind the [gui] extra.
- [X] Choose PyQt5 under the GPL, and record the split license in pyproject.toml.
- [X] Resolve the canonical-code rule: a non-adjacent repeated day is written out again.
- [X] Add name.encode and name.encode_day with their round-trip tests.
- [X] Add the dashboard subcommand to writing_habit/cli.py for parity.
- [X] Add pytest-qt to the dev extra.
- [X] Add a Makefile whose test target exports QT_QPA_PLATFORM=offscreen.
- [X] Keep a per-command database path on every command form.
- [X] Fix the canonical run rule: a change of project repeats the activity letter.
- [X] Build the window shell, the packaging, and the two graphical test modules.
- [X] Generate the command forms from argparse, and run every tracker subcommand.
- [X] Extend the same machinery to writing-schedule for the Generate and Sheets tabs.
- [X] Build the weekly-table grid in the Schedule tab, above the template and check forms.
- [X] Write the table round-trip test, which passes for all thirteen shipped tables.
- [X] Give EM and TT single-letter aliases so examples/my-week.org can be named.
- [X] Decide which project keeps the code H in examples/aug24.org, and give the other one O.
- [X] Replace :spec: in aug24.org, which now reads :risky:.
- [X] Read :risky: as the tag for the speculative class, and stop reading :speculative:.
- [X] Rewrite the tag in the nine shipped tables, the README, and tracking-formats.md.
- [X] Make the grid and the legend editable, and add the writer, Save, Save as, and the scaffold.
- [X] Port the :risky: tag to writing-habit-el, so the two ports read a table alike.
- [X] Port name.encode to writing-habit-el, with ERT tests and a cross-port check.
- [X] Add a rename action, so the canonical file name is actionable rather than advisory.
- [X] Add the preview dock for the dashboards, the plots, and the text files.
- [X] Keep PyQtWebEngine in its own preview extra, optional even with gui.
- [X] Add docs/gui.md, place it in the toctree, and illustrate it with two screenshots.
- [X] State the split license in the README and on the Sphinx installation page.
- [X] Fix tests/test_seasons.py, whose golden depended on the SQLite build underneath.
- [ ] Review the diff in both repositories, then commit; the elisp twin changed too.
- [ ] Regenerate the two GUI screenshots under assets/images with the CC BY metadata, if they should live there as well.
- [ ] Consider a writing-habit-gui section in the elisp twin's README, saying the interface is Python-only.
- [ ] Rename the two template files that merge a project change inside a generative run.
- [ ] Port name.encode to the Emacs Lisp twin so the two ports name files alike.
\end{verbatim}
\end{document}
